% 递增阈值 h_t(x)=I{x>=t} 的一记录压缩。
% 上行：8 条可实现样本，分数 0.08,0.21,0.35,0.47(标签0) 与 0.56,0.69,0.83,0.94(标签1)。
% 压缩：只保留最小正例 (0.56,1)；重构：t=0.56，判对全部8条。
% 下行：没有正例的样本 0.15,0.40,0.62,0.88 全为0，保留空集，重构为恒0函数。
% 横轴比例：分数 1 对应 100 mm。
\begingroup
\providecolor{FigureGreen}{HTML}{829C8B}
\providecolor{FigureGreenFill}{HTML}{E8EEE9}
\begin{tikzpicture}[figure picture,x=1mm,y=1mm]
% 第一行
\fill[FigureBlueFill] (56,24) rectangle (100,34);
\draw[FigureStroke,line width=.6pt] (0,26)--(100,26);
\foreach \x in {0,0.5,1}{\draw[FigureStroke] ({100*\x},25.2)--({100*\x},26.8); \node[figure note] at ({100*\x},22.5) {$\x$};}
\foreach \x in {0.08,0.21,0.35,0.47}{\draw[FigureRed,fill=white,line width=.8pt] ({100*\x},29) circle (1.3);}
\foreach \x in {0.69,0.83,0.94}{\fill[FigureBlue] ({100*\x},29) circle (1.3);}
\fill[FigureBlue] (56,29) circle (1.3);
\draw[FigureGreen,line width=1pt] (56,29) circle (2.6);
\draw[FigureStroke,line width=.9pt,dash pattern=on 2pt off 1pt] (56,24)--(56,36);
\node[figure note,anchor=south] at (56,36) {保留 $(0.56,1)$，重构 $t=0.56$};
\node[figure note,anchor=west] at (101,29) {$S$};
\node[figure note] at (78,31.8) {$h_t=1$};
\node[figure note] at (27,31.8) {$h_t=0$};
% 第二行
\draw[FigureStroke,line width=.6pt] (0,8)--(100,8);
\foreach \x in {0,0.5,1}{\draw[FigureStroke] ({100*\x},7.2)--({100*\x},8.8); \node[figure note] at ({100*\x},4.5) {$\x$};}
\foreach \x in {0.15,0.40,0.62,0.88}{\draw[FigureRed,fill=white,line width=.8pt] ({100*\x},11) circle (1.3);}
\node[figure note,anchor=west] at (101,11) {$S'$};
\node[figure note] at (50,15.5) {没有正例：保留空集，重构恒为 $0$ 的函数};
% 图例
\draw[FigureRed,fill=white,line width=.8pt] (4,41) circle (1.3); \node[figure note,anchor=west] at (6,41) {标签 $0$};
\fill[FigureBlue] (22,41) circle (1.3); \node[figure note,anchor=west] at (24,41) {标签 $1$};
\end{tikzpicture}
\endgroup
